\section*{Chapter \ref{ch:m3:defi-credit-and-leverage} --- DeFi Credit and Leverage}

\begin{solution}{exo:m3:defi-credit-and-leverage:1}
At 50\%: borrow $4\% \times 50/90 = 2.22\%$, supply $2.22\% \times 0.5 \times 0.9 = 1.00\%$. At 95\%: borrow $4\% + 60\% \times 5/10 = 34\%$,
supply $34\% \times 0.95 \times 0.9 = 29.07\%$.
\end{solution}

\begin{solution}{exo:m3:defi-credit-and-leverage:2}
$100{,}000/(50 \times 0.825) = 2{,}424.24$ dollars.
\end{solution}

\begin{solution}{exo:m3:defi-credit-and-leverage:3}
The transaction either ends with the loan and fee repaid or reverts entirely, in which case the \omterm{def:m3:blockchains-for-traders:key}{tokens} never left the pool. There is no state in
which the loan is outstanding.
\end{solution}

\begin{solution}{exo:m3:defi-credit-and-leverage:4}
$50{,}000/2{,}500 \times 1.05 = 21$ ether.
\end{solution}

\begin{solution}{exo:m3:defi-credit-and-leverage:5}
No: $(1 + \beta)\theta = 1.1 \times 0.825 = 0.9075$, above the \omterm{def:m3:defi-credit-and-leverage:ltv}{health factor} of 0.80, so each liquidation lowers it further
(\cref{prop:m3:defi-credit-and-leverage:liq}); the position is heading for bad debt.
\end{solution}

\begin{solution}{exo:m3:defi-credit-and-leverage:6}
Supply the \omterm{def:m3:defi-credit-and-leverage:staking}{liquid staking token}, borrow ether against it, stake the ether into more of the \omterm{def:m3:blockchains-for-traders:key}{token}, and repeat, earning the \omterm{def:m3:defi-credit-and-leverage:staking}{staking yield} on several
times the capital less the borrow rate on the debt. It loses if the borrow rate rises above the \omterm{def:m3:defi-credit-and-leverage:staking}{staking yield} (utilisation spikes), or if the \omterm{def:m3:blockchains-for-traders:key}{token}
trades below the ether behind it and the loops are liquidated.
\end{solution}

\begin{solution}{exo:m3:defi-credit-and-leverage:7}
At a pump of 1.93 if the \omterm{def:m3:blockchains-for-traders:key}{tokens} bought become worthless, 1.72 if they keep USD~1: borrowing 60\% of an inflated value and abandoning collateral worth
the true value pays as soon as $0.6k$ exceeds one plus the pump's cost per dollar of collateral.
\end{solution}

\begin{solution}{exo:m3:defi-credit-and-leverage:8}
Over-collateralised at the \omterm{def:m3:blockchains-for-traders:contract}{oracle}'s price. If the price can be moved, or the collateral cannot be sold near it in a crash, or liquidators do not act in
time, the collateral can be worth less than the debt: the protection is only as good as the \omterm{def:m3:blockchains-for-traders:contract}{oracle} and the liquidity of the collateral.
\end{solution}

\begin{solution}{pb:m3:defi-credit-and-leverage:1}
\textbf{1.} $0.6 \times 5 = $ USD~3 million. \textbf{2.} $k^* = 50/(0.6 \times 5) = 16.7$. \textbf{3.} $2 \times (\sqrt{16.7} - 1) = $ USD~6.16
million, buying $1 - 1/\sqrt{16.7} = 75.5\%$ of the market's \omterm{def:m3:blockchains-for-traders:key}{tokens}, 1.51 million. \textbf{4.} $50 - 5 - 6.16 = $ USD~38.84 million; USD~40.35
million if the 1.51 million \omterm{def:m3:blockchains-for-traders:key}{tokens} keep USD~1. \textbf{5.} The pool is empty: further pumping only costs. \textbf{6.} 1.93. \textbf{7.} The attack
never pays: pushing a USD~20 million market costs more than any borrowing it enables. \textbf{8.} $k^*$ doubles to 33.3, but the attack still pays
from a pump of 4.98 and can gain USD~35.4 million: halving the loan-to-value is not enough. \textbf{9.} A cap of about USD~5.8 million on borrowing
against this collateral: the gain $\min(B, 3k) - 5 - 2(\sqrt{k} - 1)$ (USD million) is then never positive. \textbf{10.} The attacker would have to hold
the pumped price for half an hour, multiplying the cost and the arbitrage losses while the price is out of line. \textbf{11.} That Eisenberg pumped MNGO
on the exchanges feeding the \omterm{def:m3:blockchains-for-traders:contract}{oracle}, raising the reported price more than thirteenfold in 30 minutes, and used his inflated positions as collateral to
withdraw over USD~110 million. \textbf{12.} On a motion under Rule 29, which asks whether the evidence could sustain the verdict, the
judge vacated the first two counts and entered a judgment of acquittal on the third (23 May 2025). \textbf{13.} The suppliers of the pool, pro rata, unless the
protocol has a reserve or a backstop. \textbf{14.} It lasts longer than a transaction and needs the attacker's own capital; a flash-loan attack borrows
the capital and must complete in one transaction, so it can only exploit prices that move within it. \textbf{15.} Freeze the collateral, lower its
loan-to-value, cap borrowing against it, change the \omterm{def:m3:blockchains-for-traders:contract}{oracle}, and vote on recoveries. \textbf{16.} Both: the design made the attack cheap, and whether the
conduct was unlawful is for courts. \textbf{17.} Only with borrow caps sized to the cost of moving its \omterm{def:m3:blockchains-for-traders:contract}{oracle}'s sources, or not at all. \textbf{18.} By the
collateral they accept and its caps, their \omterm{def:m3:blockchains-for-traders:contract}{oracles}, their bad-debt history and reserves, and utilisation (the ability to withdraw). \textbf{19.}
\emph{Named result:} a pump of 16.7 times borrows the pool dry; it costs USD~6.16 million in the source market (plus the USD~5 million of abandoned
collateral) against USD~50 million borrowed, a gain of about USD~38.8 million. \textbf{20.} The only thing it knows about the world, and so the thing an
attacker buys.
\end{solution}

\begin{solution}{iq:m3:defi-credit-and-leverage:1}
Threshold-weighted collateral over debt. Below one, anyone may repay part of the debt (up to the close factor) and take collateral at a discount, the
\omterm{def:m3:defi-credit-and-leverage:ltv}{liquidation bonus}.
\iqlookfor{the ratio, the threshold, and permissionless liquidation.}
\end{solution}

\begin{solution}{iq:m3:defi-credit-and-leverage:2}
A loan that must be repaid within the transaction or the transaction reverts. It needs atomic execution of arbitrary code across contracts, which a
\omterm{def:m3:blockchains-for-traders:chain}{blockchain} provides and traditional settlement does not.
\iqlookfor{atomicity as the collateral.}
\end{solution}

\begin{solution}{iq:m3:defi-credit-and-leverage:3}
The pool has no maturity and must keep \omterm{def:m3:blockchains-for-traders:key}{tokens} available for withdrawals; a steep rate above the target utilisation forces repayment and attracts supply
before the pool is fully lent, while a gentle slope below keeps borrowing cheap in normal times.
\iqlookfor{liquidity for withdrawals.}
\end{solution}

\begin{solution}{iq:m3:defi-credit-and-leverage:4}
From the collateral's liquidity (depth of the markets the \omterm{def:m3:blockchains-for-traders:contract}{oracle} reads and where liquidators sell), its volatility over the time liquidation takes, the
cost to manipulate its \omterm{def:m3:blockchains-for-traders:contract}{oracle}, and a borrow cap so that the maximum loss is bounded; revisit as markets change.
\iqlookfor{liquidation horizon, depth and manipulation cost.}
\end{solution}

\begin{solution}{iq:m3:defi-credit-and-leverage:5}
The yield was paid, not earned, and the peg depended on the floating \omterm{def:m3:blockchains-for-traders:key}{token}'s value exceeding the \omterm{def:m3:blockchains-for-traders:stable}{stablecoins} outstanding; when demand turned, redemptions
created floating \omterm{def:m3:blockchains-for-traders:key}{tokens} that nobody wanted. Watch the ratio of the floating \omterm{def:m3:blockchains-for-traders:key}{token}'s value to the \omterm{def:m3:blockchains-for-traders:stable}{stablecoin} supply, the source of the yield, and redemption
flows.
\iqlookfor{reflexivity and the subsidy.}
\end{solution}

\begin{solution}{iq:m3:defi-credit-and-leverage:6}
Monitor \omterm{def:m3:defi-credit-and-leverage:ltv}{health factors} from the chain's state and pending \omterm{def:m3:blockchains-for-traders:contract}{oracle} updates; simulate each liquidation with the collateral's sale path and gas; fund with a
\omterm{def:m3:defi-credit-and-leverage:flash}{flash loan} and sell in the same transaction so no inventory is carried; bid for position through bundles only up to the expected profit; and cap exposure
to collateral that cannot be sold at once.
\iqlookfor{atomic funding and disposal, and bidding discipline.}
\end{solution}
